\documentclass[11pt,a4paper]{article}
\input{../../sablony/pripravy/priprava.sty}
\usepackage{tikz}
\nastavPripravu{1. ročník}{Souhrnné opakování}{133}{Rovnice a volba metody}
\begin{document}
\blok{Úvod a připomenutí}
Před každým výpočtem se žák rozhodne: roznásobit, dosadit, sečíst, rozložit, použít diskriminant, intervaly, nebo zkoušku? V odůvodnění volby metody není nutné opisovat definice; chceme rozumět postupu. Řešení vždy uzavři množinou $K$.
\blok{Společně řešené příklady}
\begin{prikladReseniBox}[title={Závorky stačí roznásobit}]
Vyřešte $3(2x-1)-2(x+4)=x+1$.
\tcblower
\[
\begin{aligned}
6x-3-2x-8&=x+1\\
4x-11&=x+1\\
3x&=12\\
K&=\{4\}
\end{aligned}
\]

\end{prikladReseniBox}
\begin{prikladReseniBox}[title={Jedna soustava, dvě známé cesty}]
Vyřešte soustavu $\begin{cases}x+y=7,\\2x-y=2.\end{cases}$ jednou sčítáním a podruhé dosazením.
\tcblower
\[
\begin{aligned}
(x+y)+(2x-y)&=7+2\\
3x&=9\ \Rightarrow\ x=3,\ y=4\\
y&=7-x\ \Rightarrow\ 2x-(7-x)=2\\
3x&=9\ \Rightarrow\ x=3,\ y=4\\
K&=\{[3;4]\}
\end{aligned}
\]

\end{prikladReseniBox}
\begin{prikladReseniBox}[title={Kořeny hledáme vhodnou metodou}]
Vyřešte $x^2-5x+6=0$ a vysvětlete, proč je rozklad výhodnější než počítání diskriminantu.
\tcblower
\[
\begin{aligned}
x^2-5x+6&=(x-2)(x-3)\\
(x-2)(x-3)&=0\\
K&=\{2;3\}
\end{aligned}
\]

\end{prikladReseniBox}
\begin{prikladReseniBox}[title={Absolutní hodnota vyžaduje dva případy}]
Vyřešte $|2x-3|=5$.
\tcblower
\[
\begin{aligned}
2x-3&=5\quad\text{nebo}\quad2x-3=-5\\
x&=4\quad\text{nebo}\quad x=-1\\
K&=\{-1;4\}
\end{aligned}
\]

\end{prikladReseniBox}
\blok{Samostatné procvičování}
\begin{zadaniReseniBox}[title={Závorky a neznámá na obou stranách}]
Vyřešte $2(x-3)+5=3x-4$.
\tcblower
\[
\begin{aligned}
2x-1&=3x-4\\
K&=\{3\}
\end{aligned}
\]

\end{zadaniReseniBox}
\begin{zadaniReseniBox}[title={Dvě zlomkové strany}]
Vyřešte $\dfrac{x+1}{3}=\dfrac{x-2}{2}$.
\tcblower
\[
\begin{aligned}
2(x+1)&=3(x-2)\\
2x+2&=3x-6\\
K&=\{8\}
\end{aligned}
\]

\end{zadaniReseniBox}
\begin{zadaniReseniBox}[title={Soustava}]
Vyřešte $\begin{cases}2x+y=11,\\x-y=1.\end{cases}$.
\tcblower
\[
\begin{aligned}
3x&=12\\
x&=4,\quad y=3\\
K&=\{[4;3]\}
\end{aligned}
\]

\end{zadaniReseniBox}
\begin{zadaniReseniBox}[title={Parametr a výjimka}]
Diskutujte $(a-2)x=a-2$ pro $a\in\mathbb R$.
\tcblower
\[
\begin{aligned}
a\ne2&:\quad K=\{1\}\\
a=2&:\quad 0=0\ \Rightarrow\ K=\mathbb R
\end{aligned}
\]

\end{zadaniReseniBox}
\blok{Volitelná rezerva a výzva}
\poznamka{Výzva: vymysli soustavu s jedním řešením, kterou lze vyřešit sčítáním i porovnáním — a obhaj, která cesta je kratší.}
\blok{Domácí úkol}
\begin{zadaniReseniBox}[title={Domácí příprava}]
\begin{enumerate}[label=\alph*),itemsep=2pt]
\item Vyřešte $5x-7=2x+8$.
\item Vyřešte $\dfrac{2x-1}{4}-\dfrac{x+1}{2}=1$.
\item Vyřešte $\begin{cases}3x+y=14,\\x-y=2.\end{cases}$.
\item Vyřešte $x^2+x-12=0$.
\item Vyřešte $x^4-5x^2+4=0$.
\item Vyřešte $|x-2|+|x+2|=6$.
\end{enumerate}
\tcblower
\begin{enumerate}[label=\alph*),itemsep=2pt]
\item
\[
\begin{aligned}
3x&=15\\
K&=\{5\}
\end{aligned}
\]
\item
\[
\begin{aligned}
2x-1-2(x+1)&=4\\
-3&=4\quad\text{spor}\\
K&=\varnothing
\end{aligned}
\]
\item
\[
\begin{aligned}
4x&=16\\
x&=4,\quad y=2\\
K&=\{[4;2]\}
\end{aligned}
\]
\item
\[
\begin{aligned}
x^2+x-12&=(x+4)(x-3)\\
K&=\{-4;3\}
\end{aligned}
\]
\item
\[
\begin{aligned}
t&=x^2\ge0\\
(t-1)(t-4)&=0\\
t&=1\ \text{nebo}\ 4\\
K&=\{-2;-1;1;2\}
\end{aligned}
\]
\item
\[
\begin{aligned}
x<-2&:\quad-2x=6\Rightarrow x=-3\\
-2\le x\le2&:\quad4\ne6\\
x>2&:\quad2x=6\Rightarrow x=3\\
K&=\{-3;3\}
\end{aligned}
\]
\end{enumerate}
\end{zadaniReseniBox}
\poznamka{Čas: 5 + 20 + 15 + 5 minut. U delšího úkolu vyžaduj výslovně volbu metody a větu, proč jsou získaná řešení přípustná.}
\end{document}
